%% ch14 --- trace-assert w calosci: ten sam instrument w trzech jezykach
%%
%% Napisany, wrzesien 2026. Lista asercji pochodzi z agent-eval-bench,
%% skopiowana z evals/schema/scenario.schema.json oraz
%% Assertions/AssertionEvaluator.cs w commicie 12b1bbd -- patrz
%% notes/01-curriculum.md sekcja 5. Nic tutaj nie jest odtworzone z pamieci.

\chapter{trace-assert w całości: ten sam instrument w trzech językach}
\label{ch:traceassert}

\noindent
Przebieg agenta jest drogi do oceniania i tani do sprawdzenia. Ocenianie
wymaga drugiego modelu, rubryki, budżetu i argumentu za tym, że sędzia
zgadza się z człowiekiem. Sprawdzenie wymaga zapisu samego przebiegu i
funkcji, która go czyta \dash{} a większość tego, co psuje się agentom na
produkcji, widać w tym zapisie, bez wydawania przez kogokolwiek opinii o
tekście odpowiedzi.

Ten rozdział kończy \code{trace-assert}: model śladu, pełny zestaw asercji
i pakiet. Na końcu zbudujesz w Pythonie to, co w C\# umiałeś zbudować już
wcześniej, a ostatnia sekcja rozdziału to porównanie \dash{} bo ten sam
projekt istnieje w trzech językach i okazuje się, że nie niosą one tego
samego.

\section{Najtańsza ewaluacja, jaką agent może mieć}
\label{sec:ch14-cheapest}
\index{ewaluacja!deterministyczna}

Zanim przeczytasz dalej, odpowiedz sobie: czego potrzeba, żeby uruchomić
ewaluację agenta? Większość odpowiedzi wymienia model \dash{} coś musi
przeczytać odpowiedź i zdecydować, czy była dobra.

\begin{trapbox}
\textbf{Asercja nad śladem potrzebuje modelu, żeby cokolwiek ocenić.} Nie
potrzebuje, a to założenie sprawia, że większość zespołów nie ma żadnej
ewaluacji: zestaw, który potrzebuje dostawcy, potrzebuje budżetu, klucza i
sieci, więc nie uruchamia się przy pull requeście. \emph{Czy agent wywołał
narzędzie, które miał wywołać? W kolejności, w jakiej miał? Z
identyfikatorem, który wrócił z wcześniejszego wywołania? Bez wywołania
tego, czego mu zabroniono?} Każde z tych pytań jest rozstrzygalne z
zapisanego przebiegu. Dodatek~\ref{app:traps} trzyma ten wpis i to jest
powód, dla którego ta warstwa istnieje.
\end{trapbox}

\noindent
Są więc dwie warstwy i nie są konkurencją. Warstwa deterministyczna ocenia
mechanikę \dash{} co zostało wywołane, w jakiej kolejności, z czym.
Warstwa sędziowska ocenia jakość zdania, które czyta użytkownik. Ten
rozdział buduje pierwszą. To ta, która może zablokować scalenie, bo nic nie
kosztuje i dwa razy daje tę samą odpowiedź.

\begin{csbox}
Inżynier od xUnit budował to już wcześniej, nie nazywając tego ewaluacją:
to test integracyjny, który przechwytuje, co zrobił badany system, i
asercjonuje nad przechwytem, a nie nad wartością zwracaną. Nowy jest
przedmiot, nie technika \dash{} a powodem, by nadać przechwytowi nazwę i
kształt, jest to, że kilkanaście asercji nad nim chce odpowiedzi na te same
pytania.
\end{csbox}

\section{Ślad to nie log}
\label{sec:ch14-trace}
\index{ślad!model}

Oczywistym modelem przebiegu jest lista zdarzeń z rodzajem przy każdym i to
właśnie miała pierwsza wersja tego pakietu. Model ten przewraca się na
pierwszym ciekawym pytaniu. \emph{Czy zapis nastąpił przed zatwierdzeniem?}
to porównanie wywołania narzędzia ze zdarzeniem kontraktu, a jedna lista z
dyskryminatorem nie ma osi, na której dałoby się je porównać, dopóki tej
osi nie wymyślisz.

Model jest więc taki, do jakiego doszedł \code{agent-eval-bench}: trzy
równoległe zapisy, w których wywołania narzędzi i zdarzenia dzielą jeden
indeks pozycji.

\pyregion{code/src/trace_assert/trace.py}{model}{Trzy zapisy, z których składa się ślad}{lst:ch14-model}

\noindent
Warto nazwać trzy decyzje z tego listingu, bo każda z nich to defekt, który
ktoś już wypuścił na produkcję.

\textbf{Wywołanie narzędzia jest jedno, niezależnie od liczby prób
transportu.} Handler odporności, który ponawia dwa razy pod agentem, nie
wykonał trzech wywołań; pętla orkiestratora, która woła narzędzie
ponownie \dash{} owszem. To różne liczby odpowiadające na różne pytania, a
model, który je skleja, nie odróżni ponowionego żądania od podwójnej
rezerwacji.

\textbf{Wywołanie istnieje niezależnie od tego, czy się powiodło.} Gdyby
nieudany zapis nie zostawiał śladu, \emph{zapisu nigdy nie próbowano} i
\emph{zapis się nie powiódł} wyglądałyby w materiale dowodowym tak samo, a
to są dwa wyniki, które test degradacji istnieje po to, by rozróżnić.

\textbf{Nic tutaj nie pozwala asercji czytać tekstu agenta} poza
\code{Turn.reply}, które jest po to, by jedna asercja mogła poszukać w nim
wyciekłych identyfikatorów. W chwili, gdy asercja dopasowuje tekst w rodzaju
\enquote{zarezerwowałem}, każde przeredagowanie promptu staje się
czerwonym testem, a zestaw zaczyna oceniać sformułowania.

\section{Dwanaście, w trzech rodzinach}
\label{sec:ch14-twelve}
\index{ślad!asercje}

Asercji jest \val{ta.assertions} i ta liczba nie należy do tej książki:
to liczba ze schematu samego \code{agent-eval-bench}, a lista została z
niego skopiowana, nie przypomniana. Dodanie tutaj trzynastej zrobiłoby z
niej trzynastą, której dwa pozostałe porty tego samego projektu nie mają.

\mermaidfig{ta-three-families}{Asercje w liczbie \val{ta.assertions}, w trzech rodzinach, na jakie dzieli je specyfikacja}{trzy rodziny asercji}

\noindent
Obecność i nieobecność to te dwie połowy, które myli się razem. Test
dowodzący, że agent odmówił, musi asercjonować i odmowę, \emph{i} to, że
zabronione wywołanie nigdy nie nastąpiło: agent, który uprzejmie przeprasza,
a potem i tak woła narzędzie, sam z siebie przechodzi pierwszą asercję.

\pyregion{code/src/trace_assert/assertions.py}{presence}{Obecność i nieobecność, które zawsze pisze się parami}{lst:ch14-presence}

\noindent
W rodzinie kształtu siedzą te ciekawe, a \code{order} jest tą, której
specyfikację trzeba przeczytać uważnie. Nie porównuje najwcześniejszego
wystąpienia jednej rzeczy z najwcześniejszym wystąpieniem drugiej. Każde
wystąpienie drugiej musi nastąpić po \emph{pierwszym} wystąpieniu
pierwszej \dash{} bo słabsze odczytanie przepuszcza zapis wykonany przed
bramką, o ile po nim nastąpi drugi, grzeczny zapis, a bramka dotyczy
właśnie tego pierwszego.

\pyregion{code/src/trace_assert/assertions.py}{order}{\code{order} i odczytanie, które wygląda na równoważne, a nie jest}{lst:ch14-order}

\begin{exercise}{e14_01_called_with}{Przenieś jedną asercję}
\code{tool\_called\_with} pyta, czy któreś wywołanie niosło wskazane przez
ciebie argumenty, w jednym z dwóch odczytań: \code{subset} dopuszcza inne
argumenty, \code{exact} nie. Napisz ją. Test prosi o oba odczytania, o
komunikat błędu nazywający to, co zastano, i o przypadek, który odróżnia
prawdziwą asercję od ozdobnej \dash{} narzędzie, którego nigdy nie
wywołano, musi dać błąd, a nie przejść z braku kontrprzykładu.
\end{exercise}

\section{Trzy dyscypliny, z których jedna jest nośna}
\label{sec:ch14-disciplines}
\index{ewaluacja!asercja pusta}

Funkcje w liczbie \val{ta.assertions} dzielą trzy reguły. Pierwsza jest
powyżej: nic nie dopasowuje tekstu. Trzecia jest drobna \dash{} argument,
którego nie da się użyć, jest błędem, a nie przejściem, co ma tu znaczenie,
bo port w C\# ma przed sobą schemat JSON odrzucający źle napisaną asercję,
a ten port nie ma nic poza sobą.

Druga jest warta osobnej sekcji.

\begin{trapbox}
\textbf{Żadna asercja nie może przechodzić pusto.} Napisz
\code{call\_attempts} w oczywisty sposób \dash{} weź największą liczbę prób
spośród wywołań tego narzędzia, daj błąd, gdy przekracza próg \dash{} a
będzie zielona na przebiegu, w którym narzędzia w ogóle nie wywołano, bo
maksimum z niczego wynosi zero, a zero mieści się w każdym progu. Zestaw
asercjonujący powściągliwość przechodzi wtedy na przebiegu, w którym nic
się nie wydarzyło, czyli na tym jedynym, o którym powinien krzyczeć
najgłośniej.
\end{trapbox}

\pyregion{code/src/trace_assert/assertions.py}{attempts}{Większość tego to zabezpieczenie przed zaliczeniem przebiegu, w którym nic się nie wydarzyło}{lst:ch14-attempts}

\noindent
Ta sama reguła rozstrzyga trzy dalsze przypadki.
\code{order} daje błąd, gdy którejkolwiek ze stron nigdy nie było.
\code{tool\_called\_with} daje błąd, gdy narzędzia nie wywołano. A
\code{argument\_grounded} \dash{} asercja, która wyrabia całej warstwie
jej miejsce, bo łapie pewnie wymyślony identyfikator, czego żaden sędzia
nie robi niezawodnie.

\pyregion{code/src/trace_assert/assertions.py}{grounded}{Ugruntowanie jako struktura: wartość musiała wrócić z wcześniejszego wywołania}{lst:ch14-grounded}

\noindent
A komunikat błędu, który mówi tylko, która asercja padła, jest nic nie
wart, bo czerwony kolor już to powiedział. Oto, co wypisują cztery z nich
na przebiegu, który przeskoczył bramkę, wymyślił identyfikator, pokazał go
użytkownikowi i skończył, wyczerpując limit iteracji:

\transcript{ch14-broken}

\begin{exercise}{e14_02_never_vacuous}{Asercja, która może tylko przejść}
Napisz \code{call\_attempts} tak, żeby nie mogła być zielona na przebiegu,
w którym narzędzia nigdy nie wywołano. Test prosi o oba kierunki, a jednym
z nich jest przebieg pusty.
\end{exercise}

\section{Jak wygląda zestaw}
\label{sec:ch14-suite}

Rejestrator \dash{} ten z rozdziału~\ref{ch:aitoolkit} albo eksporter
czytający spany, które ustawił rozdział~\ref{ch:observability} \dash{}
spłaszcza przebieg do trzech zapisów. Wszystko dalej to funkcje nad
wartością.

\pyregion{code/ch14/eval_suite.py}{suite}{Cztery grupy asercji nad jednym zapisanym przebiegiem}{lst:ch14-suite}

\noindent
W prawdziwym zestawie każda z nich jest testem biorącym fixture
\code{trace} z rozdziału~\ref{ch:testing}, a zliczanie robi własna kolekcja
pytesta. Jako skrypt wypisuje to:

\transcript{ch14-suite}

\noindent
Własny zestaw projektu przewodniego to \val{ta.tests} testów nad asercjami
w liczbie \val{ta.assertions} i kończy się poniżej \val{ta.ceiling} sekund.
Ta liczba jest pułapem, a nie pomiarem, i to celowo: sekundy, które zajęło
to tej maszynie, są własnością tej maszyny, a liczba, której nikt inny nie
odtworzy, to liczba, przez którą kontrola dryfu świeci na czerwono przy
każdym uruchomieniu, aż ktoś ją wyłączy. Książka twierdzi tylko to, co
czytelnik może sprawdzić.

\begin{csbox}
Port w C\# zwraca z każdej asercji rekord wyniku, zamiast rzucać wyjątkiem,
i to nie jest różnica języków \dash{} to różnica w tym, co musi zrobić
harness. Ocenia on cały korpus i raportuje odsetek przejść wobec zapisanej
linii bazowej, więc potrzebuje każdego wyniku, również tych nieudanych. W
Pythonie zbiera to za ciebie framework testowy, więc asercja, która rzuca
wyjątkiem, jest kształtem krótszym i bardziej idiomatycznym.
\end{csbox}

\section{Ten sam ślad w trzech językach}
\label{sec:ch14-ports}
\index{ślad!porty}

Projekt istnieje trzy razy, a porównanie nie jest tym, które tytuł
rozdziału obiecuje bez zastrzeżeń. \textbf{W trzech językach istnieje model
śladu.} Warstwa zbudowana na nim jest w każdym inna.

\mermaidfig{ta-three-ports}{Jeden model, przeniesiony dwukrotnie, niosący w każdym języku inną warstwę}{co naprawdę niesie każdy port}

\begin{center}
\small
\begin{tabularx}{\linewidth}{@{}lXXX@{}}
\toprule
 & \textbf{.NET} & \textbf{TypeScript} & \textbf{Python} \\
\midrule
Warstwa & asercje i sędzia & tylko sędzia & tylko asercje \\
Typy & rekordy, typy referencyjne dopuszczające null & schematy \pkg{zod}, parsowane na brzegu & zamrożone dataclass, \pkg{pyright} strict \\
Async & niepotrzebny; ślad jest w pamięci & worker konsumuje kolejkę & niepotrzebny \\
Framework & xUnit v3, jeden test na scenariusz & \pkg{vitest} & \pkg{pytest}, jeden test na tezę \\
Odmowa & schemat JSON odrzuca złą asercję & schemat parsuje na brzegu & funkcja rzuca wyjątkiem \\
Pakowanie & referencja do projektu, nigdy nie publikowana & prywatny, wdrażany jako kontener & wheel w indeksie \\
\bottomrule
\end{tabularx}
\end{center}

\noindent
Z tej tabeli wynikają dwie rzeczy i żadnej z nich nie dało się przewidzieć,
zanim otwarto te trzy repozytoria.

Pierwsza jest taka, że \textbf{podróżuje kształt danych, a nie kod}. Każdy
port implementuje swoją warstwę na nowo, w idiomie swojego języka, i żadna
z implementacji nie przypomina pozostałych; identyczny jest zapis, który ta
warstwa czyta, aż do wspólnego indeksu pozycji. Projekt, który podróżuje,
to projekt, którego interfejsem jest wartość.

Druga dotyczy tego, gdzie może mieszkać kontrakt. Port w .NET odrzuca źle
napisaną asercję schematem JSON, bo jego asercje przychodzą z plików YAML
pisanych ręcznie. Port w Pythonie nie ma w obiegu żadnego pliku, więc ta
sama odmowa musi być kodem i funkcja rzuca \code{ValueError} tam, gdzie
zrobiłby to schemat. Żaden z nich nie jest staranniejszy; staranność jest
w tym samym miejscu potoku w obu przypadkach.

\begin{note}
Opis tego rozdziału, i wstęp książki, przedstawiają porównanie jako trzy
porty jednego instrumentu. To prawda o śladzie, a nie o asercjach, które
istnieją w .NET i w Pythonie. Portem w TypeScripcie jest
\code{judge-worker} i niesie on sędziego z rubryką, a nie tę warstwę; jego
własne typy śladu mówią w komentarzu, że zostały przeniesione z tego samego
pliku w C\#. Poprawiony został opis, a nie rozdział napisany pod niego.
\end{note}

\section{Publikacja}
\label{sec:ch14-publish}
\index{pakowanie!uv build}

\code{uv build} to jedno polecenie i nie ono jest warte strony. Warte
strony jest to, co z niego wyszło, bo to \dash{} a nie kod \dash{} ktoś
sobie instaluje.

\pyregion{code/ch14/publish.py}{build}{Zbuduj dystrybucję, potem odczytaj z niej metadane}{lst:ch14-build}

\noindent
Wheel to plik zip z członem metadanych w formacie nagłówków wiadomości
pocztowej, dlatego jest parsowany, a nie dzielony na dwukropkach: format
dopuszcza powtarzające się pola, a słownik zachowa tylko ostatnią z
powtórzonych linii \code{Requires-Dist}. Uruchomiony na projekcie tej
książki raportuje:

\transcript{ch14-publish}

\noindent
I to jest uczciwy wniosek tej sekcji. Katalog \code{code/} tej książki
buduje dystrybucję nazwaną po książce, przypiętą do jednej wersji minor
Pythona, na której książka jest budowana, i deklarującą szereg
zależności uruchomieniowych.
\code{trace\_assert} nie importuje żadnej z nich. To dobry projekt na
czternaście rozdziałów listingów i nie jest to pakiet, który ktokolwiek
powinien instalować. Publikuje się projekt, a projekt to jego metadane.

Sama publikacja to potem dwa polecenia i żadnego sekretu:

\begin{shellcmd}
uv build
uv publish --trusted-publishing automatic
\end{shellcmd}

\noindent
Trusted publishing oznacza, że indeks weryfikuje workflow, który wysyła
pakiet, a nie token wklejony przez kogoś w ustawieniach: nie ma
długowiecznego poświadczenia, które mogłoby wyciec, a fork nie może go
użyć. Jednorazowa konfiguracja jest po stronie indeksu, poza jakimkolwiek
plikiem w tym repozytorium, i nie została zrobiona dla własnej kopii tej
książki \dash{} to otwarte zgłoszenie, a nie krok, którego wykonanie ten
rozdział może ci pokazać.

\begin{exercise}{e14_03_publishable}{Zrób z tego pakiet do publikacji}
Napisz tabelę \code{[project]}, której potrzebuje opublikowany
\code{trace-assert}. Test ją parsuje i pyta o siedem rzeczy; przeczytaj
test, bo to on jest specyfikacją. Dwie z tych siedmiu to te, które ta
sekcja zmierzyła: pakiet nie może deklarować żadnych zależności
uruchomieniowych, bo każdy import w nim pochodzi z biblioteki standardowej,
i musi dopuszczać więcej niż tę jedną wersję minor Pythona, na której
akurat budowana jest książka. Trzecia to ta, którą obiecał
Rozdział~\ref{ch:testing}: fixture \code{trace} trafiał tam przez linijkę
\code{conftest.py}, a pakiet nie może wymagać tej linijki od swoich
użytkowników.
\end{exercise}

\section{Co masz}
\label{sec:ch14-what}

Pakiet, czyli artefakt, który ktoś może zainstalować. Projekt, który
opiszesz na rozmowie w dwóch zdaniach \dash{} \emph{ślad to trzy zapisy na
jednym indeksie, a czyta go dwanaście funkcji} \dash{} i obronisz przed
oczywistym zarzutem, czyli tym, że nie ocenia jakości i nigdy nie miał.
Oraz kawałek Pythona, który inżynier C\# rozpozna linia po linii, napisany
w idiomie tego języka, a nie jako tłumaczenie z jego własnego.

To było całe twierdzenie tej książki. Tabele odwzorowań są w
dodatku~\ref{app:cheatsheet}; nawyki, które kosztują popołudnie, w
dodatku~\ref{app:traps}; a wszystko pod \code{code/} się uruchamia.
